% 本文件由 paperswarm.tex_gate 从台账自动生成，请勿手工编辑。
% 每个数值均由证据台账注入：claim_id -> (artifact SHA-256, JSON Pointer)。
\section{Experiments}

Experiments were conducted to evaluate the performance of the minimum-norm ordinary least squares (OLS) solution and ridge regression on a synthetic dataset designed to mimic real-world challenges. The dataset consisted of 90 training samples and 200 test samples, with 120 features in the design matrix. Each configuration was averaged over 8 random seeds to ensure robustness. 

The mean test mean squared error (MSE) of the OLS solution was found to be 0.6198, with a standard deviation of 0.1329. This indicates significant variability in the OLS performance across different random splits. In contrast, ridge regression exhibited lower mean test MSE of 0.3041, with a much smaller standard deviation of 0.0623, suggesting more consistent performance. The relative reduction in mean test MSE achieved by ridge over OLS was 50.94\%, highlighting the effectiveness of ridge regression in reducing the MSE.

The selected ridge penalty (alpha) was 0.1, which was determined to be the value that minimized the mean test MSE. This penalty was crucial in balancing the trade-off between bias and variance, leading to improved generalization. Additionally, the irreducible noise floor, used as a reference for the MSE, was found to be 0.1225, indicating the inherent variability in the data that could not be reduced by any model.

These results underscore the importance of regularization in machine learning models, particularly in scenarios where the dataset might be noisy or the number of features is high. The findings also support the use of multi-agent systems like PaperSwarm for generating reproducible and transparent research, as they can systematically explore different configurations and provide reliable insights. Future work could explore the impact of different noise levels and feature distributions on the performance of these models.
